\section{Понятие цикла и особенности его отражения в экономической теории}

\subsection{Определение цикла}
Сложность и многогранность явления цикла в экономике в совокупности со существованием множества конкурирующих теоретико-методологических подходов к его анализу обуславливает отсутствие единого общепринятого определения в рамках экономической теории, в связи с чем для целей настоящего исследования является целесообразной формулировка эклектического определения, выступающего исходной абстракцией в рамках дальнейшей теоретической конкретизации данного понятия и фиксирующего признаваемые различными направлениями экономической мысли его наиболее общие моменты.

Таким образом, \textit{экономический цикл} можно определить как объективизированное в рамках закономерностей изменения объема совокупного выпуска и прочих сопутствующих показателей экономической конъюнктуры (занятости, цен, прибыли, инвестиций и так далее) волнообразное движение совокупной экономической активности, описываемое в качестве последовательной смены фаз ее качественных состояний и внешне наблюдаемое в форме устойчивых отклонений от некоторого \textquotedbl{}нормального\textquotedbl{} состояния или долгосрочного тренда развития. $\dagger$ (Приложение \ref{app:cycle-definitions}, стр. \pageref{app:cycle-definitions})

Однако данное определение носит предельно абстрактный и преимущественно описательный характер, фиксируя прежде всего внешнюю форму проявления экономического цикла и не раскрывая его фундаментального основания, механизма возникновения и специфики развития как особой формы движения обуславлвающей его экономической системы, требуя тем самым дальнейшей конкретизации посредством анализа внутренних противоречий, структурных диспропорций и закономерностей ее динамики.

\subsection{Термины цикла}
В научных трудах по экономической теории наряду с наиболее общей категорией \textit{экономического цикла} можно выделить совокупность более конкретных терминов, акцентирующих внимание на отдельных существенных сторонах, специфике внешнего проявления, механизма, продолжительности или сферы протекания цикла, вследствие чего уточнение его терминологии имеет существенное содержательно-методологическое значение, поскольку выбор того или иного термина играет основополагающую роль в контексте отражения сущности данного явления.

Исторически наиболее ранними и по смыслу относящимися к характеристике цикла выступали связанные с кризисом понятия (\textit{кризис, депрессия, затруднение, затоваривание, давление, бедствие} и др.), поскольку в действительном движении капиталистической экономики именно кризисы изначально выступали прообразом явления цикла, издавна фиксируясь различными авторами прежде всего в качестве наиболее резких, разрушительных и социально значимых проявлений внутренней неустойчивости экономической системы, в то время как теоретическое осмысление целостной и закономерной структуры цикла как такового датируется более поздним периодом. $\dagger$ (Приложение \ref{app:cycle-terms}, стр. \pageref{tab:crisis-terms})

В рамках дальнейшего становления теоретических концепций отражающие нелинейность экономического развития основные термины по своему смыслу могут быть сведены к двум основополагающим категориям:
\begin{itemize}
    \item \textit{термины цикла}, подразумевающие относительно детерминированный характер движения экономической системы, в ограниченной степени подверженный воздействию случайных отклонений (\textit{цикл, волна, ритм} и др.);
    \item \textit{термины колебаний}, подразумевающие преимущественно стохастический характер движения экономической системы, в ограниченной степени сохраняющий признаки циклической закономерности (\textit{колебания, флуктуации, осциляции, импульсы, шоки} и др.) $\dagger$ (Приложение \ref{app:cycle-terms}, стр. \pageref{tab:cycle-terms}, \pageref{tab:fluctuation-terms})
\end{itemize}

В научной литературе употребление основных терминов цикла зачастую сопровождается использованием ряда конкретизирующих понятий, в результате чего полноценные наименования экономических циклов целесообразно классифицировать по следующим критериям:
\begin{itemize}
    \item \textit{общему содержанию} (экономический / деловой / бизнес- / конъюнктурный / хозяйственный цикл и др.);
    \item \textit{уровню анализа экономики} (отраслевой цикл, цикл мировой экономики и др.);
    \item \textit{типу экономики} (цикл капиталистический / перходной / открытой экономики и др.);
    \item \textit{отрасли или сфере проявления} (ромышленный / аграрный / монетарный цикл и др.);
    \item \textit{эмпирически наблюдаемому показателю} (цикл запасов / нормы прибыли и др.);
    \item \textit{временной характеристике} (среднесрочный / долгосрочный цикл и др.);
    \item \textit{причинно-механистическому основанию} (эндогенный / экзогенный / синхронизированный цикл и др.);
    \item \textit{институциональной принадлежности} (бюджетный / политический / фискальный цикл и др.);
    \item \textit{способу теоретической идентификации} (реальный деловой цикл и др.);
    \item \textit{фамилии авторов} (цикл Жугляра / Китчина, волны Кондратьева, ритмы Кузнеца и др.) $\dagger$ (Приложение \ref{app:cycle-classification}, стр. \pageref{tab:cycle-classification-general})
\end{itemize}

Однако в контексте анализа закономерностей развития современной капиталистической экономики именно \textit{промышленный цикл} является самым подходящим по смыслу термином, отражающим закономерный, последовательный и внутренне связанный характер движения капиталистической системы и конкретизирующим его по сфере промышленного производства в качестве ее фундаментального основания, заключающего в себе динамику процессов расширения и сокращения объема совокупного выпуска, накопления и обесценения капитала в совокупности с технологической модернизацией производственных мощностей, а также кризисного разрешения накопленных противоречий и диспропорций; таким образом в наиболее полной форме систематизирующим и формализующим материальную основу, специфическую сферу проявления и внутренний механизм соответствующего цикла.

\subsection{Методология исследования цикла}
Множество подходов к исследованию цикла в экономике характеризуется существенными различиями своих методологических оснований в совокупности с пониманием природы самой экономической системы и закономерностей ее развития, в результате чего все существующие теории цикла целесообразно классифицировать по следующим категориям в зависимости от их методологии:
\begin{itemize}
    \item \textit{марксистская теория цикла}, соответствующая фундаментальным положениям марксизма как комплексного и разностороннего учения, исторически возникшего во второй половине 19 века посредством развития идей немецкой классической философии, английской политической экономии и французского утопического социализма, совмещающего в себе материалистическую диалектику, марксистскую политическую экономию и научный социализм; рассматривающая капиталистическую экономику как конкретно-историческую и внутренне противоречивую систему общественных отношений, а промышленный цикл – как специфическую форму ее диалектического развития в рамках определенного исторического этапа;
    \item \textit{буржуазные теории цикла}, соответствующие различным положениям буржуазной экономической мысли, получившей свое развитие в 17 веке начиная с классической политической экономии и впоследствии разделившейся на множество конкурирующих экономически школ, сохраняющей доминирующее положение в рамках современного экономического мейнстрима преимущественно в форме неоклассического синтеза, сочетающего в себе неоклассическое учение о рыночном равновесии, ценах и распределении с кейнсианским учением об эффективном спросе, занятости и макроэкономическом регулировании; исходящие из метафизического представления об экономике как о внутренне непротиворечивой в своей основе системе и в основном рассматривающие явление цикла через призму абсолютизации роли отдельно взятых и относительно самостоятельных факторов – денежно-кредитных, инвестиционных, технологических, психологических, институциональных и прочих;
    \item \textit{эклектические теории цикла}, занимающие промежуточное положение, в различных пропорциях соединяющие несовместимые элементы двух предыдущих категорий или внешне воспроизводящие понятийный аппарат одной из них при фактическом заимствовании содержательных положений друго.
\end{itemize}

Наиболее оптимальной методологической основой анализа цикла в экономике выступает именно \textit{марксистская теория}, рассматривающая промышленный цикл как специфическую форму движения капиталистической экономики, проистекающую из единства ее производительных сил и производственных отношений в рамках соответствующего ей способа производства, обострения и разрешения ее имманентных противоречий, динамики накопления капитала и периодического возникновения диспропорций расширенного воспроизводства, отражающую самые общие и объективно существующие тенденции развития промышленного производства в качестве ее материального основания, в максимально полной форме раскрывающую истинную сущность явления цикла, его материальную основу, внутренний механизм, исторические пределы и связь с кризисами общего перепроизводства как необходимыми моментами диалектического развития капиталистической экономики.

\subsection{Подходы к исследованию цикла}
Существующие в рамках экономической науки подходы к исследованию цикла в первую очередь целесообразно классифицировать в зависимости от их общетеоретического основания, восходящего к одной из многочисленных школ, исходящих из собственной онтологии экономической системы, по-своему определяющих ее сущность, источник и закономерности развития, степень устойчивости и характер движения, обуславливая тем самым фундаментальные расхождения между исходными предпосылками различных теорий.

В контексте сравнительного анализа теоретических подходов к исследованию цикла целесообразным является выделение нижеследующего ряда ключевых критериев:
\begin{itemize}
    \item \textit{сущность экономики} (саморегулирующаяся равновесная система обмена / денежная и макроэкономически неустойчивая система / институционально и исторически обусловленный порядок / основанная на извлечении прибавочной стоимости система капиталистического воспроизводства и др.);
    \item \textit{сущность цикла} (внешнее возмущение рыночного равновесия / следствие ошибок денежно-кредитной политики или иных экзогенных шоков / результат внутренней нестабильности инвестиций, финансов и ожиданий / объективная форма движения капитала, выражающая его внутренние противоречия и др.);
    \item \textit{сущность государственного вмешательства} (минимальное вмешательство и признание самокоррекции рынка / ограниченное или умеренное контрциклическое воздействие / активное регулирование совокупного спроса, финансовой системы и занятости и др.)
\end{itemize}

Среди всех упомянутых подходов именно \textit{марксистская политическая экономия} выступает самым последовательным и содержательным общетеоретическим основанием исследования цикла, позволяя в наиболее полной и достоверной форме анализировать капиталистическую экономику в качестве исторически определенной и внутренне противоречивой системы общественного воспроизводства; раскрыть закономерности ее развития через призму базовых законов капиталистического способа производства; объединить в рамках единого основания анализ структуры экономики, источников ее динамики, классовых отношений и исторических пределов самого капитализма; при этом диалектически рассматривая государственное регулирование как значимый фактор перераспределения, отсрочки или смягчения кризисных явлений, однако не в качестве средства устранения их коренной причины.

\subsection{Инструментарий исследования цикла} %\textcolor{red}{Черновик:}
Исследование экономического цикла не может быть локализовано в пределах какого-либо одного раздела экономической науки, поскольку сама цикличность выражает движение хозяйственной системы сразу в нескольких измерениях – как динамику совокупного производства, как нарушение и восстановление пропорций воспроизводства, как изменение отраслевой и пространственной структуры экономики, как исторически определенную форму развития капитализма и как результат воздействия политических и институциональных факторов. Поэтому проблематика цикла пронизывает как общетеоретические направления экономической мысли, так и специальные дисциплины, изучающие кризисы, конъюнктуру, накопление капитала, инвестиции, деньги, кредит, межотраслевые связи, мировое хозяйство и государственное регулирование, что делает ее по существу междисциплинарной внутри самой экономической науки.

Теоретические дисциплины, затрагивающие вопросы цикла или используемые для его исследования, целесообразно классифицировать по следующим категориям:
\begin{itemize}
    \item Экономическая динамика – теория экономических циклов, теория экономических кризисов, теория длинных волн, теория экономического роста, теория экономического развития, теория инноваций и технологического прогресса;
    \item Экономическое воспроизводство – теория воспроизводства, теория накопления капитала, теория инвестиций, теория прибыли и нормы прибыли, теория распределения доходов, теория реализации, теория денег, кредита и процента, теория основного и оборотного капитала;
    \item Координация и пропорции экономической системы – теория рыночного равновесия, теория диспропорциональности воспроизводства, теория межотраслевых балансов, теория конкуренции и межфирменного взаимодействия, теория цен и относительных цен, теория межсекторной передачи шоков;
    \item Пространственный анализ экономики – микроэкономика, макроэкономика, мезоэкономика, отраслевая экономика, международная экономика, мировая экономика, региональная и пространственная экономика;
    \item Историческое развитие и структура экономики – теория империализма, теория неравномерного развития, теория монополистического капитала, теория зависимости, мир-системный анализ, теория режимов накопления, теория финансиализации;
    \item Политико-институциональные аспекты экономики – теория экономической политики, теория фискальной политики, теория денежно-кредитной политики, теория правил и дискреции, теория общественного выбора, теория политического делового цикла, теория институциональных ограничений и реформ.
\end{itemize}

Комплексное исследование цикла требует рассмотрения всех этих аспектов в их единстве и историческом развитии, однако они не являются равноположенными: первичным выступает анализ капиталистического способа производства, отношений производства и воспроизводства, накопления капитала и присущих ему внутренних противоречий, тогда как вторичными являются специальные формы проявления этих противоречий – динамические модели, денежно-кредитные механизмы, межотраслевые диспропорции, пространственная неравномерность, институциональные формы и государственное регулирование. Иначе говоря, цикл должен исследоваться как исторически конкретная форма движения капитала, а все прочие дисциплины и разделы экономической науки должны использоваться не изолированно, а как взаимодополняющий инструментарий, позволяющий переходить от наиболее общих закономерностей капиталистического воспроизводства к их конкретным проявлениям в кризисах, фазах цикла, структурных сдвигах, территориальной дифференциации и политических механизмах регулирования. Именно такой подход позволяет избежать эклектики и встроить разнообразные аналитические средства в единую теоретическую систему, где исходным основанием служит политико-экономический анализ воспроизводства капитала, а остальные разделы конкретизируют и раскрывают отдельные стороны этого процесса.

\subsection{Методы исследования цикла} %\textcolor{red}{Черновик:}
Исследование экономического цикла опирается на широкий круг научных методов, поскольку сама цикличность представляет собой сложное и многослойное явление, сочетающее в себе внутренние закономерности воспроизводства капитала, историческую изменчивость форм кризиса, институциональные условия хозяйствования и количественно измеряемую динамику макроэкономических показателей, вследствие чего полноценный анализ цикла требует соединения различных категорий научных методов, каждый из которых раскрывает особую сторону циклического процесса.

Существующие методы исследования цикла в экономике целесообразно разделить по следующим категориям:
\begin{itemize}
    \item Философско-методологические и общенаучные методы – диалектический метод, метод единства исторического и логического, восхождение от абстрактного к конкретному, анализ и синтез, индукция и дедукция, системный подход, причинно-следственный анализ;
    \item Исторические и сравнительные методы – историко-экономический метод, историко-генетический метод, историко-сравнительный метод, метод периодизации, сравнительно-страновой и сравнительно-институциональный анализ, ретроспективный анализ кризисов, case study, метод аналогий;
    \item Политэкономические и специально-экономические методы – метод анализа общественного воспроизводства, метод анализа накопления капитала, анализ нормы прибыли, анализ инвестиционного процесса, анализ распределительных отношений, анализ денежно-кредитного механизма, институциональный, структурный и межотраслевой анализ;
    \item Качественные эмпирические и прикладные методы – экспертная оценка, опросы предприятий и потребителей, конъюнктурное наблюдение, контент-анализ деловой и финансовой информации, интервью и полевые исследования, диагностический мониторинг отраслей и рынков, качественный сценарный анализ;
    \item Статистические, математические и эконометрические методы – анализ динамических рядов, корреляционный и регрессионный анализ, панельная эконометрика, фильтры Ходрика–Прескотта, Бакстера–Кинга и Кристиано–Фитцджеральда, спектральный и вейвлет-анализ, динамические факторные модели, VAR и SVAR, модели Марковского переключения, фильтр Калмана, алгоритмы датировки поворотных точек, межотраслевые и равновесные модели;
    \item Прогностические и практико-ориентированные методы – прогнозирование по опережающим индикаторам, nowcasting, сценарное прогнозирование, стресс-тестирование, системы раннего предупреждения кризисов, машинное обучение, имитационное моделирование, форсайт-анализ, макромоделирование экономической политики.
\end{itemize}

Использование этих методов должно быть выстроено иерархически: первичными выступают философско-методологические и политэкономические методы, поскольку именно они позволяют раскрыть сущность цикла как формы движения внутренних противоречий капиталистического воспроизводства, накопления и кризиса; вторичными, но необходимыми, являются исторические, эмпирические, статистические, эконометрические и прогностические методы, которые служат для конкретизации, проверки и измерения тех закономерностей, которые уже выявлены на теоретическом уровне. Иначе говоря, исследование должно строиться по принципу восхождения от абстрактного к конкретному: от анализа наиболее общих категорий – товара, денег, капитала, прибавочной стоимости, накопления, воспроизводства – к рассмотрению конкретных форм проявления цикла в динамике инвестиций, прибыли, кредита, занятости, цен, отраслевых пропорций и исторически определенных кризисов. Такой подход позволяет избежать как абстрактного теоретизирования без эмпирической опоры, так и чисто описательного накопления фактов без раскрытия их внутренней связи, подчиняя весь исследовательский инструментарий единой задаче – выявлению сущности, механизма и конкретно-исторических форм экономического цикла.